\documentclass[11pt]{article}
\usepackage[utf8]{inputenc}
\usepackage{amsmath,amssymb}
\usepackage[margin=1in]{geometry}
\usepackage{hyperref}
\emergencystretch=3em

\title{Face moments from envelopes}
\author{Claude, with Daniel Murfet}
\date{2026-09-07}

\begin{document}
\maketitle
\sloppy

\begin{abstract}
The bundled moment hypotheses \texttt{FaceMoments} of the $d=2$ cutoff theorems are discharged from
polynomial-exponential envelopes of the Taylor data (Astra~\#5(c)). The key estimate is the finite-part
bound
\[
|\mathrm{FP}_\gamma(\Psi(\cdot,s))|\le\frac{b^{M-\gamma}}{M-\gamma}H(s)+\sum_{m<M}|f_m(s)|\,|\mathrm{axisPrim}(\gamma,b,m)|,
\]
after which every face coefficient (monomial, finite part, logarithmic, and the remainder envelope)
is bounded by a constant times $(1+s)^De^{\beta sa'}$ and the Gaussian moment lemma applies; the
exponents are automatically positive since the monomial exponent is $(m+h_1+1)/k_1$. The constructor
\texttt{faceMoments\_of\_envelope} removes the opaque \texttt{FaceMoments} assumptions from
applications. Zero \texttt{sorry}/\texttt{axiom}.
\end{abstract}

\section{The things formalised}
{\small
\begin{verbatim}
theorem axisFinitePart_abs_le (γ b H : ℝ) (hb : 0 < b) (M : ℕ) (hM : γ < M) (f : ℝ → ℝ)
    (fm : ℕ → ℝ) (hrem : ∀ v ∈ Ioc (0 : ℝ) b, |taylorRem f fm M v| ≤ H * v ^ M) :
    |axisFinitePart γ b f fm M|
      ≤ H / ((M : ℝ) - γ) * b ^ ((M : ℝ) - γ)
        + ∑ m ∈ Finset.range M, |fm m| * |axisPrim γ b m|
noncomputable def faceScalar (γ A : ℝ) (k₁ k₂ : ℕ) (m : ℕ) : ℝ
theorem faceCoeff_eq_scalar_mul ... :
    faceCoeff γ A k₁ k₂ f m s = faceScalar γ A k₁ k₂ m * f m s
theorem faceExp_pos (p₂ : ℝ) (h₁ k₁ : ℕ) (hk₁ : 0 < k₁) (m : ℕ) :
    0 < faceExp p₂ ((k₁ : ℝ) * p₂ - h₁ - 1) k₁ m
theorem faceLogCoeff_abs_le ... :
    |faceLogCoeff γ k₁ k₂ f M s| ≤ ∑ m ∈ Finset.range M, 1 / ((k₁ : ℝ) * k₂) * |f m s|
theorem faceMoments_of_envelope (β p₂ b a' Cf CH : ℝ) (h₁ k₁ k₂ M D : ℕ) (hβ : 0 < β)
    (hb : 0 < b) (hk₁ : 0 < k₁) (hk₂ : 0 < k₂) (hp₂ : 0 < p₂) (Ψ : ℝ → ℝ → ℝ)
    (hΨc : Continuous (Function.uncurry Ψ)) (f : ℕ → ℝ → ℝ) (hf : ∀ m, Measurable (f m))
    (H : ℝ → ℝ) (hH : Measurable H) (hH0 : ∀ s, 0 ≤ H s) (hγ : (k₁ : ℝ) * p₂ - h₁ - 1 < M)
    (hCH : 0 ≤ CH)
    (hfenv : ∀ m, m < M → ∀ s, 0 < s → |f m s| ≤ Cf * (1 + s) ^ D * Real.exp (β * s * a'))
    (hHenv : ∀ s, 0 < s → H s ≤ CH * (1 + s) ^ D * Real.exp (β * s * a'))
    (hrem : ∀ s, 0 < s → ∀ u ∈ Ioc (0 : ℝ) b,
      |Ψ u s - ∑ m ∈ Finset.range M, f m s * u ^ m| ≤ H s * u ^ M) :
    FaceMoments β p₂ b h₁ k₁ k₂ M Ψ f H
\end{verbatim}
}

\section{Mathematics}

The regularised integral is the tail estimate of unit~100 at $\varepsilon=b$:
$|\int_0^bv^{-1-\gamma}R_M|\le H\,b^{M-\gamma}/(M-\gamma)$; the endpoint counterterms are bounded
termwise. Each face coefficient is then dominated by a constant times the common envelope
$(1+s)^De^{\beta sa'}$: monomial coefficients are scalar multiples of $f_m$
(\texttt{faceScalar}), the logarithmic coefficient is a sum of at most $M$ such multiples, the
finite part by the displayed bound, and the remainder envelope is a constant times $H$. The
Gaussian moment lemma (\texttt{moment\_integrableOn\_of\_envelope}) needs only the exponent
$\alpha-1>-1$: the monomial exponents are $(m+h_1+1)/k_1>0$, the finite-part and log exponents are
$p_2>0$, and the remainder exponent $p_2+(M-\gamma)/k_1>p_2$. Because the weight is Gaussian, any
finite $a'$ is admissible (Astra~\#5(c)).

\section{Paper-facing meaning}

With this constructor the cutoff theorems of units~107 and~109 need only: continuity of the data,
compatible Taylor remainders of the face jets with envelopes $C(1+s)^De^{\beta sa'}$, and the
mixed remainder bound. For the chart amplitude $\eta e^{\beta s\xi}$ these are exactly the outputs of
the smooth rectangular Taylor theory (Astra~\#5 ranks~3 and~5), the next units.

\section{Lean notes}

\begin{itemize}
\item \texttt{weighted\_taylorRem\_tail\_le} at $\varepsilon=b$ gives the regularised-integral bound
for free.
\item Field goals generated by \texttt{refine ⟨?\_, ?\_, ?\_⟩} for a \texttt{Prop} structure with
\texttt{Fin M ⊕ Fin 2} indices: \texttt{rcases i with m | i} then \texttt{fin\_cases i}, with
\texttt{simpa [faceα, faceJ, faceC] using …} to align the unfolded index functions.
\item A \texttt{set} variable and its unfolded value are defeq but not syntactically equal; goals
created after the \texttt{set} (structure fields) show the unfolded form, and \texttt{rw} still
closes them by \texttt{rfl}.
\item \texttt{abs\_of\_nonneg (by positivity : (0 : ℝ) ≤ 1 / (↑k₁ * ↑k₂))} works without
positivity hypotheses on the naturals.
\end{itemize}

\end{document}
